\documentclass[9pt]{extarticle}

\usepackage[T1]{fontenc}
\usepackage{lmodern}
\usepackage[letterpaper,top=0.40in,bottom=0.40in,left=0.46in,right=0.46in]{geometry}
\usepackage{microtype}
\usepackage[table]{xcolor}
\usepackage{booktabs}
\usepackage{tabularx}
\usepackage{array}
\usepackage{enumitem}
\usepackage{multicol}
\usepackage{titlesec}
\usepackage[hidelinks]{hyperref}

\definecolor{rulegray}{HTML}{9AA0A6}
\definecolor{rowgray}{HTML}{F3F4F6}
\setlength{\parindent}{0pt}
\setlength{\parskip}{1.8pt}
\setlength{\columnsep}{0.25in}
\setlength{\columnseprule}{0.2pt}
\def\columnseprulecolor{\color{rulegray}}
\setlist[itemize]{leftmargin=1.1em,itemsep=0.8pt,topsep=1pt,parsep=0pt}
\setlist[enumerate]{leftmargin=1.25em,itemsep=1.1pt,topsep=1pt,parsep=0pt}
\titleformat{\section}{\bfseries\small}{}{0pt}{}
\titlespacing*{\section}{0pt}{3.2pt}{1pt}
\renewcommand{\arraystretch}{1.05}
\newcolumntype{Y}{>{\raggedright\arraybackslash}X}
\newcommand{\model}[1]{\texttt{#1}}

\begin{document}
\pagestyle{empty}

\begin{center}
  {\fontsize{14}{15.5}\selectfont\bfseries
  Has Speaker Diarization Actually Improved?\\[-1pt]
  Evaluating Five Years of Model Progress Beyond Standard Adult Speech}

  \vspace{1.5pt}
  {\itshape Page 1: paper concept for external evaluation}
\end{center}

\section*{Framing and aim}

Speaker diarization has moved from embedding--clustering pipelines to
overlap-aware neural systems, end-to-end speaker-activity models, and
generative speech models that emit speaker labels and timestamps. Improvements
on adult benchmarks do not establish whether this progress transfers to less
represented conversational settings. This study evaluates four architectural
families under one protocol on conventional adult meetings, African-accented
medical conversations, adult--child interaction, and English-primary
code-switched conversation. It asks \textbf{whether architectural progress
produces gains across domains, whether those gains close transfer gaps, and
which diarization failures have or have not improved}.

The task is anonymous within-recording speaker attribution---\emph{who spoke
when}---not identification across recordings. ASR is not an outcome: when a
generative system returns words, only its speaker labels and timing are scored.

\begin{multicols}{2}
\raggedcolumns

\section*{Final evaluation datasets}

{\fontsize{7.0}{8.0}\selectfont
\begin{tabularx}{\linewidth}{@{}>{\bfseries}p{0.24\linewidth}p{0.27\linewidth}Y@{}}
\toprule
Dataset & Final subset & Role and coverage \\
\midrule
AMI Mix-Headset & 16 sessions; 9.06 h; one 3-speaker and fifteen 4-speaker
meetings & Conventional adult reference; 63.78 min overlap. \\
\rowcolor{rowgray}
AfriSpeech-Dialog & Medical only: 17 sessions; 1.77 h; all 2-speaker &
African-accented medical transfer; overlap insufficient. \\
Playlogue v1 & Official test: 34 sessions; 8.11 h; all 2-speaker & Adult--child
transfer; 2.25 h child speech (40.71\%); released RTTM has no overlap. \\
\rowcolor{rowgray}
Bangor Miami & 28 English-primary sessions; 14.85 h; 26 two-, one three-, and
one four-speaker & Code-switch transfer; 98.48 min overlap. \\
\bottomrule
\end{tabularx}}

The suite contains \textbf{95 sessions and 33.79 h}. Across the three target
datasets, 77 of 79 sessions (97.5\%) are dyadic, reducing speaker-count
confounding. Playlogue is retained as the closest speaker-count match to
AfriSpeech and Bangor; child and adult errors will be reported separately.

\section*{Research questions and measures}

{\fontsize{7.0}{8.0}\selectfont
\begin{tabularx}{\linewidth}{@{}>{\bfseries}p{0.28\linewidth}Y@{}}
\toprule
Question & Purpose and primary measures \\
\midrule
RQ1. Have successive architectural shifts improved accuracy across domains? &
Group comparison using DER, JER, miss, false alarm, confusion, count accuracy,
planned family contrasts, relative error reduction, and confidence intervals. \\
\rowcolor{rowgray}
RQ2. Have these gains closed gaps for unconventional and less represented
domains? & Transfer gaps relative to AMI, relative gap reduction,
family-by-dataset interaction, rank correlation, and rank reversals. \\
RQ3. In which speech conditions have gains and transfer been most effective? &
Errors by overlap, boundary, short turn, counting, and confusion; boundary
precision/recall/F1 and family-by-condition interaction. \\
\bottomrule
\end{tabularx}}

RQ1 tests overall progress; RQ2 tests transfer rather than raw difficulty; RQ3
locates improvements and persistent failures. Overlap claims are restricted to
AMI and Bangor. Playlogue supports child/adult, boundary, short-turn, counting,
and confusion analyses.

\columnbreak

\section*{Architecture taxonomy}

The strata follow the review taxonomy of
\href{https://doi.org/10.1016/j.csl.2021.101317}{Park \emph{et al.} (2022)}
and the clustering/EEND/separation comparison of
\href{https://doi.org/10.1016/j.csl.2023.101534}{Serafini \emph{et al.}
(2023)}. We operationalize their architectural distinctions as G1--G3 and add
G4 as an explicit post-review generative extension. The G labels are study
strata, not a claimed consensus chronology.

{\fontsize{7.0}{8.0}\selectfont
\begin{tabularx}{\linewidth}{@{}>{\bfseries}p{0.24\linewidth}Y@{}}
\toprule
Family & Defining shift and representative candidate \\
\midrule
G1: Embedding--clustering cascade & Explicit VAD, embeddings, and global
clustering. \emph{NeMo MarbleNet + TitaNet-Large + spectral clustering}. \\
\rowcolor{rowgray}
G2: Neuralized overlap-aware modular & Learned powerset/multiscale activity
with clustering for global identity. \emph{pyannote Community-1}. \\
G3: End-to-end discriminative & Direct frame-level multispeaker activity
without a clustering backend. \emph{NVIDIA Streaming Sortformer 4-speaker}. \\
\rowcolor{rowgray}
G4: Unified generative & Time-aligned anonymous speaker tokens, often generated
with words. \emph{VibeVoice-ASR; eligibility to be confirmed in the pilot}. \\
\bottomrule
\end{tabularx}}

Final selection requires a public versionable checkpoint, reproducible
inference, anonymous time-aligned output, mixed-speaker compatibility, and a
feasible common protocol. Download counts are at most an adoption tie-breaker.
Moshi remains related work because its user/system streams are predefined.

\section*{Experimental protocol}

All models receive the same 16-kHz mono recordings and versioned manifest. The
primary condition is fully automatic; an oracle-speaker-count condition
separates counting from assignment. Outputs are converted to \textbf{Rich
Transcription Time Marked (RTTM)} format. Complete recordings are used when
supported; any required chunking and cross-window speaker stitching will be
fixed after the model pilot and reported as a constraint.

Primary scoring uses zero collar with overlap included where annotated; a
0.25-second collar tests boundary sensitivity. Results are macro-averaged by
recording and use recording-clustered bootstrap intervals. Analyses control or
stratify for speaker count, overlap, turn duration/rate, speech proportion,
speaker imbalance, signal quality, and dataset. Claims concern transfer and
observed gaps, not isolated causal effects of accent, age, or domain.

\section*{Expected contribution and novelty}

The contribution is a \textbf{cross-family, architecture-aware transfer
audit}, not a new model: (1) a common protocol spanning four families and four
settings; (2) evidence of whether newer systems narrow gaps in
African-accented medical, adult--child, and English-primary code-switched
speech; (3) links between architectural shifts and overlap, boundary,
short-turn, counting, and confusion errors; and (4) a reproducible distinction
between modular, end-to-end, and generative diarization.

\end{multicols}

\newpage

\begin{center}
  {\fontsize{14}{15.5}\selectfont\bfseries Dataset Progress and
  Experiment-to-Paper Roadmap}

  \vspace{1.5pt}
  {\itshape Page 2: completed work, remaining experiments, and completion gates}
\end{center}

\section*{Collection and preparation status}

{\fontsize{7.25}{8.25}\selectfont
\begin{tabularx}{\linewidth}{@{}>{\bfseries}p{0.12\linewidth}p{0.43\linewidth}Y@{}}
\toprule
Dataset & Completed & Remaining reference action \\
\midrule
AMI & Sixteen official Mix-Headset test recordings, RTTM, and Un-partitioned
Evaluation Map (UEM) files standardized and validated. & Freeze versions and
checksums; confirm comparison-paper scoring conventions. \\
\rowcolor{rowgray}
AfriSpeech & Forty-six usable conversations prepared; final view selects the 17
medical conversations. & Export medical-only manifest; manually audit filtered
nonpositive timestamp pairs; do not infer roles from anonymous speaker IDs. \\
Playlogue & Thirty-four official test recordings trimmed, resampled, and paired
with official RTTM/UEM. & Preserve official split; document zero overlap in all
158 released RTTMs; validate child/adult-conditioned scoring. \\
\rowcolor{rowgray}
Bangor & Twenty-eight English-primary recordings converted from timestamped
CHAT to RTTM/UEM; code-switch metadata retained. & Align 35 untimed trusted
tiers or mask surrounding intervals; retain all sessions. \\
\bottomrule
\end{tabularx}}

All prepared audio is 16-kHz mono. Path, duration, RTTM/UEM, and boundary
validation completed with zero structural errors. The prepared superset has
124 recordings and 38.81 h; the locked medical-only evaluation view will have
\textbf{95 recordings and 33.79 h}. Source audio remains unchanged and the
pipeline is reproducible with \model{prepare\_inference\_datasets.py}.

\begin{multicols}{2}
\raggedcolumns

\section*{Experiments still required}

\begin{enumerate}
  \item \textbf{Lock references:} resolve Bangor tiers, export medical-only
  AfriSpeech, add checksums/version IDs, and freeze metric eligibility.
  \item \textbf{Finalize models:} verify one candidate per family; freeze
  checkpoints, licenses, settings, context limits, compute, and parsers.
  \item \textbf{Balanced pilot:} sample every dataset and speaker-count regime;
  test parsing, overlap support, count behavior, determinism, runtime, memory,
  cost, and full-recording versus window policy.
  \item \textbf{Primary inference:} run every model on every recording with
  automatic count; retain raw output, RTTM, logs, failures, runtime, memory, and
  configuration hashes.
  \item \textbf{Controlled inference:} repeat with oracle speaker count where
  supported; test the preregistered chunk/stitch condition only if required.
  \item \textbf{Overall scoring:} DER, JER, miss, false alarm, confusion,
  count error, malformed-output rate, and failure rate under zero and
  0.25-second collars.
  \item \textbf{RQ3 scoring:} short turns, boundaries, rapid changes,
  imbalance, and per-speaker confusion; overlap on AMI/Bangor; child/adult on
  Playlogue; switched/non-switched turns on Bangor.
  \item \textbf{Statistics:} recording-level macro averages and clustered
  bootstrap intervals; family, dataset, family-by-dataset, and
  family-by-condition tests; corrected planned contrasts, relative gap
  reduction, rank correlation, and rank reversal.
  \item \textbf{Robustness audit:} both collars; with/without incidental or
  highly imbalanced speakers; stratified manual error inspection; parser versus
  acoustic-model failures.
\end{enumerate}

\columnbreak

\section*{Final-paper analyses and artifacts}

\begin{itemize}
  \item \textbf{Table 1:} research questions, purposes, and measures.
  \item \textbf{Table 2:} families, defining shifts, selected checkpoints,
  parameter/access status, and output capabilities.
  \item \textbf{Dataset table:} sessions, hours, speaker counts, child/adult
  share, code-switch prevalence, overlap, turn structure, and limitations.
  \item \textbf{Main results:} DER/JER and components by family $\times$
  dataset, confidence intervals, and transfer gaps.
  \item \textbf{Primary figure:} architectural-family trajectory on AMI versus
  each target dataset.
  \item \textbf{Diagnostic figure:} gain heatmap for overlap, short turns,
  boundaries, counting, and confusion; non-eligible cells marked, not imputed.
  \item \textbf{Reproducibility package:} locked manifests, lawful data
  preparation, RTTM/UEM, parsers, scoring configuration, model versions,
  environment lockfile, seeds, and raw-to-result provenance.
\end{itemize}

\section*{Completion gates}

Full paper writing begins after:

\begin{itemize}
  \item \textbf{G1---Data:} references and metric eligibility are frozen.
  \item \textbf{G2---Pilot:} all four parsers pass the balanced pilot.
  \item \textbf{G3---Inference:} the model--recording matrix is complete or
  failures are explicitly documented.
  \item \textbf{G4---Analysis:} RQ1--RQ3 tables, intervals, and sensitivity
  analyses are generated.
\end{itemize}

Write in this order: Methods and dataset table; Results; Discussion and
limitations; Related Work and taxonomy justification; then Abstract and
Introduction after the central empirical claim is known.

\section*{Immediate next step}

Finalize the four model representatives and run the balanced pilot. The
dataset suite is otherwise mechanically ready for inference, subject to the
small reference-quality actions listed above.

\end{multicols}

\end{document}
